\pdfminorversion=4
\pdfobjcompresslevel=0
\documentclass[ignorenonframetext,]{beamer}
\setbeamertemplate{caption}[numbered]
\setbeamertemplate{caption label separator}{:}
\setbeamercolor{caption name}{fg=normal text.fg}
\usepackage{amssymb,amsmath}
\usepackage{colortbl}
\usepackage{booktabs}
\usepackage{pifont}
\usepackage{graphicx}
\usepackage{ifxetex,ifluatex}
\usepackage{fixltx2e} % provides \textsubscript
\usepackage{lmodern}
\usepackage{pdfpages}
%\setbeamercolor{background canvas}{bg=}
\makeatletter\beamer@ignorenonframefalse\makeatother
% see https://tex.stackexchange.com/questions/96581/incorrect-frame-numbers-when-pages-included-using-pdfpages-in-beamer
\def\getpdfpages#1#2{\begingroup
	\setbeamercolor{background canvas}{bg=}
	\includepdf[pages={#1}, trim=0mm 1mm 0mm 0mm,%
	pagecommand={\global\setcounter{framenumber}{%
			\the\numexpr\value{framenumber}+1\relax}%
		\expandafter\def\expandafter\insertshorttitle\expandafter{%
			\insertshorttitle\hfill\insertframenumber\,/\,\inserttotalframenumber}}%
	]{#2}
	\endgroup}

\ifxetex
  \usepackage{fontspec,xltxtra,xunicode}
  \defaultfontfeatures{Mapping=tex-text,Scale=MatchLowercase}
  \newcommand{\euro}{€}
\else
  \ifluatex
    \usepackage{fontspec}
    \defaultfontfeatures{Mapping=tex-text,Scale=MatchLowercase}
    \newcommand{\euro}{€}
  \else
    \usepackage[T1]{fontenc}
    \usepackage[utf8]{inputenc}
      \fi
\fi
% use upquote if available, for straight quotes in verbatim environments
\IfFileExists{upquote.sty}{\usepackage{upquote}}{}
% use microtype if available
\IfFileExists{microtype.sty}{\usepackage{microtype}}{}
\usepackage{url}

\providecommand{\tightlist}{%
  \setlength{\itemsep}{0pt}\setlength{\parskip}{0pt}}

% Comment these out if you don't want a slide with just the
% part/section/subsection/subsubsection title:
\AtBeginPart{
  \let\insertpartnumber\relax
  \let\partname\relax
  \frame{\partpage}
}
\AtBeginSection{
  \let\insertsectionnumber\relax
  \let\sectionname\relax
  \frame{\sectionpage}
}
\AtBeginSubsection{
  \let\insertsubsectionnumber\relax
  \let\subsectionname\relax
  \frame{\subsectionpage}
}

\setlength{\parindent}{0pt}
\setlength{\parskip}{6pt plus 2pt minus 1pt}
\setlength{\emergencystretch}{3em}  % prevent overfull lines
\setcounter{secnumdepth}{0}
\usepackage[english]{babel}
\usepackage{dsfont}
\usepackage{verbatim}
\usepackage{amsmath}
\usepackage{amssymb}
\usepackage{amsthm}
\usepackage{amsfonts}
\usepackage{csquotes}
\usepackage{multirow}
\usepackage{longtable}
\usepackage{enumerate}
% \usepackage[absolute,overlay]{textpos}
\usepackage{psfrag}
\usepackage{algorithm}
\usepackage{algpseudocode}
\usepackage{eqnarray}
\usepackage{arydshln}
\usepackage{tabularx}
\usepackage{placeins}
\usepackage{tikz}
% see https://tex.stackexchange.com/questions/32918/how-can-i-decrease-spaces-between-equations
%\usepackage{setspace}
\usepackage[nodisplayskipstretch]{setspace}
\setstretch{1.0}
\usetikzlibrary{shapes,arrows,automata,positioning,calc}
\usepackage{subfig}
% \usepackage{paralist}
\usepackage{graphicx}
\usepackage{array}
\usepackage{framed}
\usepackage{excludeonly}
\usepackage{fancyvrb}
\usecolortheme{dove}
% \usefonttheme{serif}
\usepackage{xfrac}
\usepackage{xcolor}
\usepackage{mdframed}
\usepackage{caption}
\captionsetup[figure]{labelformat=empty}
\usepackage{transparent}
\usepackage{bm}

% URL color:
\definecolor{myorange}{RGB}{80,224,176}
% Theme color of metropolis
\definecolor{metropolis_theme_color}{RGB}{38,38,38}
% Upper and lower rules of code chunks
\definecolor{darkblue}{RGB}{38,38,38}
% Font color of footer:
\definecolor{mygrey}{RGB}{240,240,240}

\definecolor{accents}{RGB}{80,224,176}

\usepackage{hyperref}
\hypersetup{
    colorlinks = true,
    linkcolor = {black},
    urlcolor = {accents},
    linkbordercolor = {white}
}


%% %%%%%%%%%%%%%%%%%%%%%%%%%%%%%%%%%%%%%%%%%%%%%%%%%%%%%%%%%%%%%%%
%% METROPOLIS THEME + CUSTOMIZATIONS
%% %%%%%%%%%%%%%%%%%%%%%%%%%%%%%%%%%%%%%%%%%%%%%%%%%%%%%%%%%%%%%%%

%\usetheme{metropolis}
%\usetheme{metropolis}
%\useoutertheme{metropolis}
%\useinnertheme{metropolis}
\setbeamercolor{background canvas}{bg=white}
\AtBeginSection{
  \let\insertsectionnumber\relax
  \let\sectionname\relax
    \begin{frame}[c,plain,noframenumbering]
    \sectionpage
    \small{\tableofcontents[currentsection]}
    \end{frame}
}

% remove space of center enviroment which often occurs around includegraphics produced from knitr plots
% https://tex.stackexchange.com/questions/74212/how-can-i-change-the-whitespace-above-and-below-center
\let\oldcenter\center
\let\oldendcenter\endcenter
\renewenvironment{center}{\setlength\topsep{0pt}\oldcenter}{\oldendcenter}

% reduce margins / border
% \setbeamersize{text margin left=10pt,text margin right=10pt}

% Reduce space between code chunks and code output in rmarkdown beamer presentation: https://stackoverflow.com/questions/35734525/reduce-space-between-code-chunks-and-code-output-in-rmarkdown-beamer-presentatio
%\setlength{\parskip}{0pt}
%\setlength{\OuterFrameSep}{-20pt}
\makeatletter
\preto{\@verbatim}{\parskip=0pt \topsep=0pt \partopsep=0pt}
\makeatother

%% Add LMU template:
%% ---------------------------------------------------------------
%\usepackage{./../../_setup/lmu-lecture}
\usepackage{../../style/lmu-lecture}
\input{../../style/common.tex}
% \usefonttheme{structurebold}
% \RequirePackage[T1]{fontenc}
% \RequirePackage[scaled=0.92]{helvet}   %% Helvetica for sans serif
% \defbeamertemplate*{frametitle}{lecture}[1][left]
% {
%   \ifbeamercolorempty[bg]{frametitle}{}{\nointerlineskip}%
%   \@tempdima=\textwidth%
%   \advance\@tempdima by\beamer@leftmargin%
%   \advance\@tempdima by\beamer@rightmargin%
%   \begin{beamercolorbox}[sep=0.2cm,#1,wd=\the\@tempdima]{frametitle}
%     \if@tempswa\else\csname beamer@fte#1\endcsname\fi%
%     {\usebeamerfont{frametitle}\rule[-0.5ex]{0pt}{2.3ex}\insertframetitle\par}%
%     \if@tempswa\else\vskip-.2cm\fi% set inside beamercolorbox... evil here...
%   \end{beamercolorbox}%
% }
% \def\beamer@fteright{\vskip0.35cm\advance\leftskip by 1.7cm\advance\rightskip by1.7cm}
% \setbeamertemplate{frametitle continuation}[from second][{\small/~\insertcontinuationcount}]
% % ------------------------------------------------------------------------
% % Geometry
% \setbeamersize{text margin left=0.8cm,text margin right=0.8cm}
% \pgfdeclarehorizontalshading{footlineshade}{4mm}{%
%   color(0pt)=(black);%
%   color(1.0\paperwidth)=(structure!50!black)}
% % redefine \ref (it has been redefined somewhere by the beamerclass)
% \def\lectureref#1{\expandafter\@setref\csname r@#1\endcsname\@firstoftwo{#1}}
% % counter for framenumber for current lecture
% \newcounter{lectureframenumber}
% % adjust framenumbers for lecture (check whether reference is already defined)
% \def\addjustlectureframenumber#1{\ifx#1\relax\else%
%   \addtocounter{lectureframenumber}{-\lectureref{lect:@@\thelecture}}\fi}

% % ------------------------------------------------------------------------
% % Navigation symbols
% \setbeamertemplate{navigation symbols}{}
% % ------------------------------------------------------------------------
% % No head lines
% \defbeamertemplate*{headline}{lecture theme}{}
% \setbeamertemplate{frametitle}{\expandafter\uppercase\expandafter\insertframetitle}








%% Add footer:
%% ---------------------------------------------------------------
\setbeamertemplate{footline}[text line]{%
    \noindent\hspace*{\dimexpr-\oddsidemargin-1in\relax}%
     \colorbox{white}{
     \makebox[\dimexpr\paperwidth-2\fboxsep\relax]{
     \color{black}
     \begin{minipage}[c][4.5ex][c]{0.5\linewidth}
       \secname
     \end{minipage}
     \hfill\begin{minipage}[c][4.5ex][c]{0.5\linewidth}
       \flushright
       \insertframenumber{}~/~\inserttotalframenumber~~
     \end{minipage}
     }}%
  \hspace*{-\paperwidth}
}

%% Add titlegraphic:
%% ---------------------------------------------------------------
% \titlegraphic{
%   \transparent{0.7}
%   \includegraphics[width=5cm]{./../../_setup/images/logos/eds_logo_black.png}
% }

%% %%%%%%%%%%%%%%%%%%%%%%%%%%%%%%%%%%%%%%%%%%%%%%%%%%%%%%%%%%%%%%%
%% NICER CODE BACKGROUND
%% %%%%%%%%%%%%%%%%%%%%%%%%%%%%%%%%%%%%%%%%%%%%%%%%%%%%%%%%%%%%%%%

% Define Shaded if not defined already through knitr:
%\definecolor{shadecolor}{RGB}{228,228,228}
\makeatletter
\@ifundefined{Shaded}{%
\newenvironment{Shaded}{\begin{snugshade}}{\end{snugshade}}%
}{}
\makeatother

\renewenvironment{Shaded}{
  %\vspace{5pt}
  \begin{mdframed}[backgroundcolor=black!8,
                   linecolor=darkblue,
                   rightline=false,
				   leftline=false]}{
  \end{mdframed}}


%% %%%%%%%%%%%%%%%%%%%%%%%%%%%%%%%%%%%%%%%%%%%%%%%%%%%%%%%%%%%%%%%
%% CHANGE URL COLOR
%% %%%%%%%%%%%%%%%%%%%%%%%%%%%%%%%%%%%%%%%%%%%%%%%%%%%%%%%%%%%%%%%

\renewcommand\UrlFont{\color{accents}}
% \hypersetup{
%     colorlinks=true,
%     linkcolor=myorange
% }

%% %%%%%%%%%%%%%%%%%%%%%%%%%%%%%%%%%%%%%%%%%%%%%%%%%%%%%%%%%%%%%%%
%% INCLUDE LATEX MATH + OTHER CUSTOMIZATIONS
%% %%%%%%%%%%%%%%%%%%%%%%%%%%%%%%%%%%%%%%%%%%%%%%%%%%%%%%%%%%%%%%%


% Include latex-math
\DeclareOldFontCommand{\sf}{\normalfont\sffamily}{\mathsf}

% \input{./../../latex-math/basic-math.tex}
% \input{./../../latex-math/basic-ml.tex}
% \input{./../../latex-math/ml-bagging.tex}
% \input{./../../latex-math/ml-boosting.tex}
% \input{./../../latex-math/ml-gp.tex}
% \input{./../../latex-math/ml-mbo.tex}
% \input{./../../latex-math/ml-nn.tex}
% \input{./../../latex-math/ml-svm.tex}
% \input{./../../latex-math/ml-trees.tex}
% \InputIfFileExists{./../../latex-math/ml-lm.tex}{}{}
% \InputIfFileExists{./../../latex-math/ml-automl.tex}{}{}
% Additional commands:
\newcommand{\abs}[1]{\ensuremath{\left| #1 \right|}}
%\renewcommand{\pikN}{\hat{\pi}^\Np} % This collides with z



% basic latex stuff
\newcommand{\pkg}[1]{{\fontseries{b}\selectfont #1}} %fontstyle for R packages
\newcommand{\lz}{\vspace{0.5cm}} %vertical space

\newcommand{\oneliner}[1] % Oneliner for important statements
{\begin{block}{}\begin{center}\begin{Large}#1\end{Large}\end{center}\end{block}}

\let\code=\texttt
\let\proglang=\textsf

\setkeys{Gin}{width=0.9\textwidth}

%% %%%%%%%%%%%%%%%%%%%%%%%%%%%%%%%%%%%%%%%%%%%%%%%%%%%%%%%%%%%%%%%
%% EXERCISE COUNTER
%% %%%%%%%%%%%%%%%%%%%%%%%%%%%%%%%%%%%%%%%%%%%%%%%%%%%%%%%%%%%%%%%

% Command for an automatic exercise counter for the exercise sheets:
\newcounter{Exercise}
\newcommand{\exercise}{
  \stepcounter{Exercise}
  \ifnum \theExercise = 1
    \Large{\textbf{Exercise \theExercise}} \\
  \else
    \vspace{\baselineskip} \Large{\textbf{Exercise \theExercise}} \\
  \fi
}

% Command for an automatic exercise counter for the courses:
\newcounter{CourseExercise}
\newcommand{\oneExercise}{
  \stepcounter{CourseExercise}
    \begin{center}
      \huge{Exercise \Roman{CourseExercise}}
    \end{center}
}
\newcommand{\twoExercises}{
  \stepcounter{CourseExercise}
    \begin{center}
      \huge{Exercise \Roman{CourseExercise} \& \stepcounter{CourseExercise} \Roman{CourseExercise}}
    \end{center}
}
\newcommand{\threeExercises}{
  \stepcounter{CourseExercise}
    \begin{center}
      \huge{Exercise \Roman{CourseExercise}, \stepcounter{CourseExercise} \Roman{CourseExercise} \& \stepcounter{CourseExercise} \Roman{CourseExercise}}
    \end{center}
}

\newcommand{\resetExercises}{\setcounter{CourseExercise}{0}}

\usetikzlibrary{shapes,arrows,automata,positioning,calc,chains,trees, shadows}
\tikzset{
  %Define standard arrow tip
  >=stealth',
  %Define style for boxes
  punkt/.style={
    rectangle,
    rounded corners,
    draw=black, very thick,
    text width=6.5em,
    minimum height=2em,
    text centered},
  % Define arrow style
  pil/.style={
    ->,
    thick,
    shorten <=2pt,
    shorten >=2pt,}
}
\usepackage{subfig}
%\beamerdefaultoverlayspecification{<+->}

\title{Best Practices for Competition Success}
\date{}

\begin{document}

\section{Best Practices for Competition Success}
%TODO graphics with overview over steps
\begin{frame}{Step 1: Understand the Problem Statement}
\begin{itemize}
  \item \textbf{Clarify Objectives \& Business Impact}: Determine whether it’s a classification or regression task, the domain context, and what “success” looks like.
  \item \textbf{Target Variable}: Understand how the target is defined and distributed.
  \item \textbf{Constraints \& Permissions}: External data limits, privacy, or regulatory issues.
  \item \textbf{Action Points}:
    \begin{itemize}
      \item Annotate key requirements and assumptions.
      \item Research the domain for feature ideas.
      \item Restate the problem in your own words to ensure clarity.
    \end{itemize}
\end{itemize}
\end{frame}

\begin{frame}{Step 2: Study the Evaluation Metric}
\begin{itemize}
  \item \textbf{Metric Matters}: Each competition or project has a specific metric (e.g., F1-score, AUC-ROC, RMSE), which guides how you tune and compare models.
  \item \textbf{Classification vs. Regression Metrics}:
    \begin{itemize}
      \item Accuracy, F1-score, AUC-ROC, and Log Loss for classification.
      \item RMSE, MAE, R\textsuperscript{2} for regression.
    \end{itemize}
  \item \textbf{Optimization Direction}: Decide if you must maximize or minimize the metric; some metrics (e.g., AUC-ROC) may need a proxy loss for training.
  \item \textbf{Action Points}:
    \begin{itemize}
      \item Simulate prediction changes to see impact on the metric.
      \item If allowed, implement custom losses aligned to the competition metric.
    \end{itemize}
\end{itemize}
\end{frame}

\begin{frame}{Step 3: Perform an Initial Data Analysis (EDA)}
\begin{itemize}
  \item \textbf{Data Profiling}: Summarize shape, basic statistics, feature and target distribution.
  \item \textbf{Quality Checks}: Identify missing values, outliers, or suspicious patterns; confirm no hidden data leakage.
  \item \textbf{Train-Test Shift}: Check if train and test set differ significantly in their feature distributions.
  \item \textbf{Action Points}:
    \begin{itemize}
      \item Use tools like Pandas Profiling or Sweetviz for automated EDA.
      \item Document anomalies or domain-specific oddities.
      \item Examine correlations, scatter plots, and histograms.
    \end{itemize}
\end{itemize}
\end{frame}

\begin{frame}{Step 4: Develop a Baseline Model}
\begin{itemize}
  \item \textbf{Why a Baseline?} A simple model (e.g., Logistic Regression, Decision Tree) gives you a reference performance.
  \item \textbf{Pipeline Check}: Make sure that data loading, preprocessing, and validation splits are correct.
  \item \textbf{Diagnostics} : Analyze misclassifications/ residuals to spot improvement areas.
  \item \textbf{Action Points}:
    \begin{itemize}
      \item Start with a naive or default-parameter model.
      \item Use cross-validation to assess reliability.
      \item Record baseline metrics to quantify future gains.
    \end{itemize}
\end{itemize}
\end{frame}

\begin{frame}{Step 5: Set Up a Robust Validation Strategy}
\begin{itemize}
  \item \textbf{Importance of Validation}: Avoid overfitting or a single random split.
  \item \textbf{Common Methods}:
    \begin{itemize}
      \item K-Fold (balanced data).
      \item Stratified K-Fold (imbalanced classes).
      \item Time-Based or Group Splits (time series, repeated measures).
    \end{itemize}
  \item \textbf{Action Points}:
    \begin{itemize}
      \item Match validation approach to the final test scenario.
      \item Keep a holdout set for final checks.
      \item Use consistent folds across experiments to compare fairly.
    \end{itemize}
\end{itemize}
\end{frame}

\begin{frame}{Step 6: Feature Engineering and Data Preprocessing}
\begin{itemize}
  \item \textbf{Missing Data}: Mean/median imputation, advanced methods, or flags.
  \item \textbf{Encoding Categorical Variables}: One-hot/ target/ frequency enc. (leakage?!).
  \item \textbf{Scaling \& Transformations}: Normalize or log-transform skewed features, especially for sensitive models (SVMs, neural nets).
  \item \textbf{New Features}: Time-based, domain knowledge, interactions (e.g. ratios).
  \item \textbf{Action Points}:
    \begin{itemize}
      \item Add features in small batches and track validation changes.
      \item Use scikit-learn Pipelines to avoid leakage.
      \item Consider dimensionality reduction if feature space is large.
    \end{itemize}
\end{itemize}
\end{frame}

\begin{frame}{Step 7: Select and Tune Models Strategically}
\begin{itemize}
  \item \textbf{Model Choices}:
    \begin{itemize}
      \item Tree-based ensembles (XGBoost, LightGBM, CatBoost) for tabular data.
      \item Neural networks for large-scale or unstructured data.
      \item Simple linear/logistic models can still compete with good feature engineering.
    \end{itemize}
  \item \textbf{Hyperparameter Tuning}:
    \begin{itemize}
      \item Grid Search, Random Search, Bayesian Optimization.
      \item Early stopping and regularization (L1, L2) to avoid overfitting.
    \end{itemize}
  \item \textbf{Action Points}:
    \begin{itemize}
      \item Track experiments (MLflow, Weights \& Biases, or spreadsheets).
      \item Start with defaults and refine gradually.
      \item Use IML tools (SHAP, permutation importance) for feature insights.
    \end{itemize}
\end{itemize}
\end{frame}

\begin{frame}{Step 8: Track Experiments and Iterate}
\begin{itemize}
  \item \textbf{Experiment Logging}: Keep detailed records of data versions, hyperparams, code revisions.
  \item \textbf{Error Analysis}: Study misclassifications and residuals to guide feature tweaks.
  \item \textbf{Iterative Improvements}:
    \begin{itemize}
      \item Refine features, adjust validation, revisit transformations.
      \item Maintain a “changelog” for each iteration.
    \end{itemize}
  \item \textbf{Action Points}:
    \begin{itemize}
      \item Automate repeated tasks (submissions, data prep).
      \item Prune ineffective ideas swiftly.
      \item Ensure reproducibility (version control, environment snapshots).
    \end{itemize}
\end{itemize}
\end{frame}

\begin{frame}{Step 9: Understand the Leaderboard and Avoid Overfitting}
\begin{itemize}
  \item \textbf{Public vs. Private Leaderboard}: Public boards use only part of the test set—overfitting leads to dramatic rank changes later.
  \item \textbf{Submission Strategy}:
    \begin{itemize}
      \item Rely mainly on internal validation.
      \item Keep a private holdout to confirm final performance.
    \end{itemize}
  \item \textbf{Leaderboard Shakeups}: Watch for big changes when private scores are revealed; it often indicates overfitting to the public set.
  \item \textbf{Action Points}:
    \begin{itemize}
      \item Limit minor tweaks solely to improve a public rank.
      \item If a jump seems suspicious, re-check data handling for leakage.
    \end{itemize}
\end{itemize}
\end{frame}

\begin{frame}{Step 10: Collaborate and Learn from Others}
\begin{itemize}
  \item \textbf{Study Top Solutions}: Many winning approaches are documented in blog posts or Kaggle discussions.
  \item \textbf{Teamwork}:
    \begin{itemize}
      \item Collaboration brings diverse perspectives and model ensembling.
      \item Share code and features for synergy.
    \end{itemize}
  \item \textbf{Community \& Networking}:
    \begin{itemize}
      \item Discussion forums reveal dataset quirks, best practices.
      \item Local meetups or online communities keep you updated.
    \end{itemize}
  \item \textbf{Action Points}:
    \begin{itemize}
      \item Document and share your findings.
      \item Adapt others’ ideas carefully—validate them on your own splits.
      \item Engage with peers for new tools and techniques.
    \end{itemize}
\end{itemize}
\end{frame}

\begin{frame}{Final Advice}
\begin{itemize}
  \item \textbf{Start Simple, Then Iterate}: Don’t jump into complex models, get a solid baseline.
  \item \textbf{Prioritize Validation}: A robust validation strategy prevents nasty surprises.
  \item \textbf{Log Everything}: Document transformations, parameters, and model versions.
  \item \textbf{Aim for Generalizable Solutions}: Don’t just chase leaderboard scores—focus on reproducibility and reliability.
\end{itemize}
\end{frame}



\end{document}
